% =====================================================================
%  SIMULAZIONE DI MATEMATICA - 1SP -- IIS "David Maria Turoldo", Zogno (BG)
%  Classe: exam  |  Compilazione: pdflatex (due passaggi)
%  Soluzioni:        decommentare \printanswers
%  Versione BES/DSA: impostare \versioneBEStrue
%  Struttura:        \iniziofila{A} \begin{esercizi} ... \end{esercizi} \finefila
%                    \iniziofila{B} \begin{esercizi} ... \end{esercizi} \finefila
%  Ogni esercizio:   \begin{esercizio}
	%                      \obiettivo{...}
	%                      \question[punti]  (oppure \question + parts con \part[punti])
	%                      testo ...
	%                    \end{esercizio}
%  Nessuno spazio per le risposte: gli studenti scrivono sul foglio a protocollo.
% =====================================================================
\documentclass[11pt,a4paper,addpoints]{exam}

% ---------- Lingua e codifica ----------
\usepackage[utf8]{inputenc}
\usepackage[T1]{fontenc}
\usepackage[italian]{babel}
\usepackage{lmodern}
\usepackage{microtype}

% ---------- Matematica, grafica, tabelle ----------
\usepackage{amsmath,amssymb}
\usepackage{siunitx}
\sisetup{output-decimal-marker={,}, per-mode=symbol, group-separator={\,}}
\usepackage{tikz}
\usetikzlibrary{arrows.meta,calc,patterns,angles,quotes}
\usepackage{pgfplots}
\pgfplotsset{compat=1.18}
\usepackage{graphicx}
\usepackage{multicol}
\usepackage{array,tabularx,booktabs}
\usepackage{xcolor}

% ---------- Pagina ----------
\usepackage[a4paper,top=2cm,bottom=2.2cm,left=2cm,right=2.2cm,
headheight=30pt,headsep=12pt]{geometry}

% ---------- Opzioni ----------
%\printanswers              % Decommentare per mostrare le soluzioni.
\ifdefined\CORRETTORE\printanswers\fi % Opzione per compilare il correttore.
\newif\ifversioneBES
\versioneBESfalse              % <-- \versioneBEStrue per la versione BES/DSA
\newif\iftabellabonus
\tabellabonusfalse

% ---------- Interlinea ----------
%  \interlinea      : testo normale
%  \interlineaparts : sottopunti a), b), ... (piu' ariosa, separa le equazioni)
\newcommand{\interlinea}{1.05}
\newcommand{\interlineaparts}{1.05}
\ifversioneBES
\usepackage[scaled=1.05]{helvet}
\renewcommand{\familydefault}{\sfdefault}
\renewcommand{\interlinea}{1.5}
\renewcommand{\interlineaparts}{1.7}
\fi
%\linespread{1.5}

% ---------- Dati della prova (da compilare) ----------
\newcommand{\materia}{Matematica}
\newcommand{\classe}{1SP}
\newcommand{\indirizzo}{Liceo Scientifico a indirizzo sportivo}
%   Liceo Scientifico | Liceo Scientifico a indirizzo sportivo | Liceo delle Scienze Umane
\newcommand{\tipoprova}{Simulazione di verifica - Raccolta di allenamento}
\newcommand{\argomento}{Numeri interi e razionali}
\newcommand{\durata}{vedi indicazioni}
\newcommand{\annoscolastico}{2026/2027}
\newcommand{\docente}{Prof. Rubens Longhi}
\newcommand{\fila}{A}
\newcommand{\indicazioni}{Svolgi sul foglio a protocollo. Mostra i passaggi essenziali e riduci i risultati. Nelle divisioni tra monomi considera le lettere al denominatore diverse da zero. \textbf{È consentito l'uso della calcolatrice.}\ifversioneBES\ Sono consentiti gli strumenti compensativi previsti dal PDP.\fi}
% ---------- Intestazione e pie' di pagina ----------
%  Ogni fila ricomincia da pagina 1: la sua prima pagina ha l'intestazione
%  della scuola e la numerazione "Pagina x di y" e' calcolata sulla singola fila.
\pagestyle{headandfoot}
\firstpageheader
{\scriptsize\textsc{IIS ``David Maria Turoldo''}\\[-1pt]
	\scriptsize Via Ronco 11 -- 24019 Zogno (BG)}
{}
{\scriptsize A.S.\ \annoscolastico\\[-1pt]\scriptsize \docente}
\firstpageheadrule
\runningheader{\small\materia\ -- \classe}{\small\tipoprova}{\small Fila \fila}
\runningheadrule
\newcommand{\paginefila}{\pageref{fine@fila@\fila}}
\firstpagefooter{}{\small Pagina \thepage\ di \paginefila}{}
\runningfooter{}{\small Pagina \thepage\ di \paginefila}{}

% ---------- Formato degli esercizi ----------
%  Riga 1:  "1."  in grassetto + obiettivo in grigio ........ (n punti)
%  Riga 2+: testo dell'esercizio
% ---------- Formato degli esercizi ----------
\newcommand{\obiettivocorrente}{}
\newcommand{\obiettivo}[1]{%
	\gdef\obiettivocorrente{#1}}
\newcommand{\stampaobiettivo}{%
	\ifx\obiettivocorrente\empty\else
	{\footnotesize\color{black!55}\textit{Obiettivo: \obiettivocorrente}}%
	\fi}
\qformat{\textbf{\thequestion.}\enspace\stampaobiettivo\hfill\colonnapunti{(\totalpoints\ \points)}\gdef\obiettivocorrente{}}
\bonusqformat{\textbf{\thequestion.}\enspace{\footnotesize\color{black!55}\textit{(facoltativo)}}\enspace\stampaobiettivo\hfill\colonnapunti{(+\totalbonuspoints\ bonus)}\gdef\obiettivocorrente{}}
\renewcommand{\questionlabel}{}
\renewcommand{\partlabel}{\thequestion\thepartno)}
\renewcommand{\subpartlabel}{\thesubpart.}
\renewcommand{\choicelabel}{\thechoice)}
\pointpoints{punto}{punti}
\bonuspointpoints{punto bonus}{punti bonus}
\newcommand{\colonnapunti}[1]{{\small #1}}
\pointformat{\small(\thepoints)}
\bonuspointformat{\small(+\thepoints)}

% Sottopunti con interlinea maggiore (le soluzioni tornano all'interlinea normale)
\renewcommand{\partshook}{\linespread{\interlineaparts}\selectfont\setlength{\itemsep}{3pt}\setlength{\rightmargin}{1cm}}

% Numerazione degli esercizi: ogni fila riparte da 1 (esercizio 1, 2, ...);
% il contatore interno di exam prosegue da una fila all'altra, cosi' punti,
% totali e tabelle restano corretti. Usare l'ambiente "esercizi" al posto di "questions".
\newcounter{offsetfila}
\renewcommand{\thequestion}{\number\numexpr\value{question}-\value{offsetfila}\relax}
\newenvironment{esercizi}
{\begin{questions}\setcounter{question}{\value{offsetfila}}}
	{\end{questions}}

% ---------- Esercizi non spezzati tra due pagine ----------
%  \begin{esercizio} ... \end{esercizio} racchiude \obiettivo + \question + parts.
%  Se non ci sta, l'esercizio passa intero alla pagina successiva.
%  Nella versione con soluzioni l'esercizio puo' spezzarsi (i correttori sono lunghi).
\cancelspace   % le soluzioni nascoste non lasciano spazio ne' punti di interruzione
\newenvironment{esercizio}
{\par\ifprintanswers\else\begin{samepage}\fi}
	{\ifprintanswers\else\end{samepage}\fi\par\vspace{3pt}}

% ---------- Soluzioni e scelte multiple ----------
\renewcommand{\solutiontitle}{\noindent\textbf{Soluzione.}\enspace}
\unframedsolutions
\SolutionEmphasis{\color{blue!60!black}\linespread{\interlinea}\selectfont}
\CorrectChoiceEmphasis{\color{blue!60!black}\bfseries}
\checkboxchar{$\square$}
\checkedchar{$\blacksquare$}

% ---------- Tabella punteggi ----------
\hqword{\textbf{Esercizio}}
\hpword{\textbf{Punti massimi}}
\hsword{\textbf{Punti ottenuti}}
\htword{\textbf{Totale}}
\chqword{\textbf{Esercizio}}
\chpword{\textbf{Punti massimi}}
\chbpword{\textbf{Bonus}}
\chsword{\textbf{Punti ottenuti}}
\chtword{\textbf{Totale}}

% ---------- Macro utili ----------
\newcommand{\R}{\mathbb{R}}
\newcommand{\N}{\mathbb{N}}
\newcommand{\Z}{\mathbb{Z}}
\newcommand{\Q}{\mathbb{Q}}
\newcommand{\dd}{\,\mathrm{d}}

% ---------- Inizio e fine di ogni fila ----------
%  \iniziofila{A} ... \finefila   poi   \iniziofila{B} ... \finefila
%  \iniziofila: nuova pagina 1, intestazione scuola, titolo, nome, indicazioni.
%  \finefila:   tabella punteggi della sola fila, voto, firma.
\newcommand{\iniziofila}[1]{%
	\clearpage
	\renewcommand{\fila}{#1}%
	\setcounter{page}{1}%
	\setcounter{offsetfila}{\value{totalq@fila}}%
	\begin{center}
		{\small\bfseries \tipoprova\ifprintanswers\ -- Soluzioni e criteri di correzione\fi}\\[3pt]
		{\Large \argomento}\\[2pt]
		{\small \indirizzo\ -- Classe \classe\ -- Durata: \durata\ -- \textbf{Fila #1}}
	\end{center}
	\vspace{4pt}
	\noindent
	\begin{tabularx}{\textwidth}{@{}l X l p{3.2cm}@{}}
		\textbf{Nome e Cognome:} & \hrulefill & \textbf{Data:} & \hrulefill \\[8pt]
	\end{tabularx}
	\vspace{2pt}
	\noindent\fbox{\parbox{\dimexpr\textwidth-2\fboxsep-2\fboxrule}{\small\linespread{1.2}\selectfont
			\textbf{Indicazioni.} \indicazioni\ Punteggio totale: \pointsinrange{fila#1} punti.}}
	\vspace{2pt}
	\begingradingrange{fila#1}}
\newcounter{totalq@fila}
\newcommand{\finefila}{%
	\setcounter{totalq@fila}{\value{question}}%
	\endgradingrange{fila\fila}%
	\vfill
	\begin{center}\small
		\iftabellabonus
		\partialcombinedgradetable{fila\fila}[h][questions]%
		\else
		\partialgradetable{fila\fila}[h][questions]%
		\fi
	\end{center}
	\vspace{6pt}
	\noindent{\small\textbf{Voto:} \rule{3cm}{0.4pt}\hfill
		\textbf{Firma del docente:} \rule{5cm}{0.4pt}}%
	\label{fine@fila@\fila}}


% Adattamenti per la simulazione: nessun punteggio, soluzioni attivabili.
\renewcommand{\colonnapunti}[1]{}
\qformat{\textbf{\thequestion.}\enspace\stampaobiettivo\hfill\gdef\obiettivocorrente{}}
\runningheader{\small\materia\ -- \classe}{\small Simulazione}{\small $\Z$, $\Q$ e applicazioni}
\renewcommand{\iniziofila}[1]{%
\clearpage\renewcommand{\fila}{#1}\setcounter{page}{1}
\begin{center}
{\small\bfseries \tipoprova\ifprintanswers\ -- Soluzioni\fi}\\[3pt]
{\Large \argomento}\\[2pt]
{\small Valore assoluto, M.C.D. e m.c.m., proporzionalità e percentuali}\\[2pt]
{\small \indirizzo\ -- Classe \classe}
\end{center}
\noindent\textbf{Nome e cognome:} \hrulefill\quad\textbf{Data:} \rule{2.3cm}{0.4pt}\par\medskip
\noindent\fbox{\parbox{\dimexpr\textwidth-2\fboxsep-2\fboxrule}{\small
\textbf{Indicazioni.} La raccolta è volutamente più ampia di una verifica di 60 minuti.
Svolgi sul foglio a protocollo, mostrando i passaggi e riducendo le frazioni ai minimi termini.
Lavora senza calcolatrice. Per ogni vero/falso, giustifica le affermazioni vere con un ragionamento generale e confuta quelle false con un controesempio.
\textbf{Allenamento in 60 minuti:} svolgi 1a, 1b, 1e, 2c, 3a, 5a, 7a, 9a, 10a, 11b, 14a, 16, 18a e 22a.}}
\medskip}
\renewcommand{\finefila}{\label{fine@fila@\fila}}
\newcommand{\sezione}[1]{\par\medskip\noindent{\large\bfseries #1}\par\smallskip}
\begin{document}
\linespread{\interlinea}\selectfont
\iniziofila{S}
\ifprintanswers\small\renewcommand{\partshook}{\setlength{\itemsep}{0pt}\setlength{\rightmargin}{1cm}}\fi
\sezione{Ragionare sulle proprietà}
\begin{esercizi}
\begin{esercizio}
\obiettivo{distinguere somme nulle e prodotti nulli}
\question Stabilisci se ciascuna affermazione è vera o falsa.
\begin{parts}
\part Se $a,b\in\Z$ e $a+b=0$, allora $a=0$ e $b=0$.
\begin{solution}
\textbf{Falsa.} $a=3$, $b=-3$: la somma è zero, ma nessuno dei due numeri è zero.
\end{solution}
\part Se $a,b\in\Z$ e $|a|+|b|=0$, allora $a=0$ e $b=0$.
\begin{solution}
\textbf{Vera.} I due valori assoluti sono non negativi; se uno fosse positivo, anche la somma sarebbe positiva. Quindi entrambi sono zero e perciò $a=b=0$.
\end{solution}
\part Se $a,b\in\Q$ e $a^2+b^2=0$, allora $a=0$ e $b=0$.
\begin{solution}
\textbf{Vera.} I quadrati sono non negativi. Una loro somma nulla impone che siano entrambi nulli; il quadrato di un numero è zero soltanto quando il numero è zero.
\end{solution}
\part Se $a,b\in\Z$ e $a^2+b^2=0$, allora $a=0$ e $b=0$.
\begin{solution}
\textbf{Vera.} Vale lo stesso ragionamento del punto precedente, anche per gli interi.
\end{solution}
\part Se $a,b\in\Z$ e $|a|\cdot|b|=0$, allora $a=0$ e $b=0$.
\begin{solution}
\textbf{Falsa.} $a=0$, $b=5$. La legge di annullamento del prodotto impone che \emph{almeno uno} dei fattori sia nullo, non che lo siano entrambi.
\end{solution}
\part Se $a,b\in\Q$ e $ab=0$, allora almeno uno tra $a$ e $b$ è zero.
\begin{solution}
\textbf{Vera.} È la legge di annullamento del prodotto. Se entrambi fossero diversi da zero, anche il prodotto sarebbe diverso da zero.
\end{solution}
\end{parts}
\end{esercizio}

\begin{esercizio}
\ifprintanswers\newpage\fi\obiettivo{argomentare su valore assoluto e numeri razionali}
\question Stabilisci se ciascuna affermazione è vera o falsa.
\begin{parts}
\part Per ogni $a\in\Z$, il numero $-a$ è negativo.
\begin{solution}
\textbf{Falsa.} Con $a=-4$ si ha $-a=4>0$.
\end{solution}
\part Se $a,b\in\Q$ e $|a|=|b|$, allora $a=b$.
\begin{solution}
\textbf{Falsa.} $a=2$, $b=-2$: stesso valore assoluto ma numeri diversi.
\end{solution}
\part Per ogni $a,b\in\Q$ vale $|a+b|=|a|+|b|$.
\begin{solution}
\textbf{Falsa.} $a=2$, $b=-2$: a sinistra $0$, a destra $4$.
\end{solution}
\part Per ogni $a,b\in\Q$ vale $|ab|=|a|\cdot|b|$.
\begin{solution}
\textbf{Vera.} Il modulo del prodotto è il prodotto delle grandezze dei fattori: con segni concordi il prodotto è non negativo, con segni discordi è negativo ma il modulo ne cambia il segno; con un fattore nullo entrambi i membri sono zero.
\end{solution}
\part Fra due numeri razionali distinti esiste sempre un altro numero razionale.
\begin{solution}
\textbf{Vera.} Se $a<b$, il numero $(a+b)/2$ è razionale ed è strettamente compreso fra $a$ e $b$.
\end{solution}
\part Moltiplicando una frazione propria positiva per una frazione impropria positiva si ottiene sempre un numero maggiore di $1$.
\begin{solution}
\textbf{Falsa.} $\frac12\cdot\frac32=\frac34<1$.
\end{solution}
\end{parts}
\end{esercizio}

\newpage\sezione{Interi, segni e valore assoluto}

\begin{esercizio}
\obiettivo{confrontare numeri e valori assoluti}
\question Inserisci $<$, $=$ o $>$ e mostra il calcolo quando necessario.
\begin{parts}
\part $|-8|\ \ldots\ |-11|$;\quad $-|-7|\ \ldots\ -6$;\quad $|-5+2|\ \ldots\ |-5|+|2|$.
\begin{solution}
$8<11$; $-7<-6$; $3<7$.
\end{solution}
\part $|18-|-10||\ \ldots\ \operatorname{mcm}(2,8)$;\quad $\operatorname{MCD}(12,18)\ \ldots\ |-3|$;\quad $-3^2\ \ldots\ (-3)^2$.
\begin{solution}
$8=8$; $6>3$; $-9<9$.
\end{solution}
\end{parts}
\end{esercizio}

\begin{esercizio}
\obiettivo{interpretare opposto, modulo e distanza}
\question Rispondi motivando.
\begin{parts}
\part Sia $a$ un intero non nullo. Quando $-a$ è positivo? Quando $-|a|$ è positivo? Il numero $|-a|$ può essere negativo?
\begin{solution}
$-a>0$ se $a<0$; $-|a|$ non è mai positivo (è strettamente negativo); $|-a|$ non può essere negativo, anzi è positivo perché $a\ne0$.
\end{solution}
\part Siano $a<0$, $b>0$ e $|a|<|b|$. Confronta $a$ con $-b$, poi $-a$ con $b$. Quanto distano sulla retta i punti $-7$ e $4$? Scrivi un'espressione con il valore assoluto.
\begin{solution}
Da $|a|=-a$ e $|b|=b$ segue $-a<b$, quindi $a>-b$. La distanza è $|4-(-7)|=11$.
\end{solution}
\end{parts}
\end{esercizio}

\begin{esercizio}
\obiettivo{calcolare espressioni in interi}
\question Calcola mostrando i passaggi.
\begin{parts}
\part $\displaystyle -18:\bigl[(-3)^2-12\bigr]+(-2)^3-|-7+3|\cdot(-2)$.
\begin{solution}
La parentesi è $9-12=-3$ e il modulo è $4$. Quindi $6-8+8=\boxed{6}$.
\end{solution}
\part $\displaystyle \bigl[(-2)^5:(-2)^3-3^2\bigr]^2-2\cdot|-8+3|-(-3)^3$.
\begin{solution}
$[4-9]^2-2\cdot5+27=25-10+27=\boxed{42}$.
\end{solution}
\end{parts}
\end{esercizio}

\begin{esercizio}
\obiettivo{usare le proprietà delle potenze}
\question Semplifica applicando le proprietà delle potenze; esplicita almeno una proprietà usata.
\begin{parts}
\part $\displaystyle \frac{(-3)^4\cdot(-3)^3}{(-3)^5}-\bigl[(-2)^3\bigr]^2:2^4$.
\begin{solution}
Nel primo termine $(-3)^{4+3-5}=9$; nel secondo $(-2)^6:2^4=2^2=4$. Risultato $\boxed5$. Si usano prodotto e quoziente di potenze con la stessa base e potenza di potenza.
\end{solution}
\part $\displaystyle \bigl[(-5)^6:(-5)^4\bigr]:5^2+(-1)^{37}-(-1)^{24}+7^0$.
\begin{solution}
$25:25-1-1+1=\boxed0$. Gli esponenti pari danno una potenza positiva, quelli dispari conservano il segno della base; $7^0=1$.
\end{solution}
\end{parts}
\end{esercizio}

\newpage\sezione{Rappresentazioni e calcolo in $\Q$}

\begin{esercizio}
\obiettivo{passare fra frazioni, decimali e percentuali}
\question Trasforma in frazioni ridotte ai minimi termini e ordina in senso crescente.
\begin{parts}
\part $-1{,}25;\quad -\frac76;\quad 0{,}\overline3;\quad 35\%;\quad \frac38$.
\begin{solution}
$-\frac54<-\frac76<\frac13<\frac7{20}<\frac38$. Si può usare il denominatore comune $120$: $-150,-140,40,42,45$.
\end{solution}
\part $-0{,}\overline6;\quad -65\%;\quad 0{,}1\overline6;\quad \frac15;\quad 0{,}2\overline3$.
\begin{solution}
$-\frac23<-\frac{13}{20}<\frac16<\frac15<\frac7{30}$. A denominatore $60$: $-40,-39,10,12,14$.
\end{solution}
\end{parts}
\end{esercizio}

\begin{esercizio}
\obiettivo{usare equivalenza e scrittura decimale}
\question Completa e giustifica.
\begin{parts}
\part Riduci $\frac{84}{126}$, $\frac{45}{120}$, $\frac{42}{175}$. Stabilisci per ciascuna se la scrittura decimale è finita, periodica semplice o periodica mista.
\begin{solution}
$\frac23=0{,}\overline6$ (semplice); $\frac38=0{,}375$ (finita); $\frac6{25}=0{,}24$ (finita). Conta il denominatore della frazione \emph{ridotta}: soli fattori $2$ e $5$ danno una scrittura finita.
\end{solution}
\part Completa $\frac{-18}{24}=\frac{\square}{4}=\frac{15}{\square}$ e $\frac7{12}=\frac{\square}{60}$. Enuncia la proprietà invariantiva. Classifica poi la scrittura decimale di $\frac7{12}$.
\begin{solution}
I numeri mancanti sono $-3,-20,35$. Moltiplicando o dividendo numeratore e denominatore per uno stesso numero diverso da zero si ottiene una frazione equivalente. $\frac7{12}=0{,}58\overline3$: periodica mista.
\end{solution}
\end{parts}
\end{esercizio}

\begin{esercizio}
\obiettivo{calcolare espressioni con frazioni}
\question Calcola e riduci il risultato ai minimi termini.
\begin{parts}
\part $\displaystyle \left[\frac34-\left(\frac56-\frac13\right)\right]:\left(-\frac12\right)+\left(-\frac23\right)^2$.
\begin{solution}
La parentesi tonda vale $\frac12$, la quadra $\frac14$. Quindi $-\frac12+\frac49=\boxed{-\frac1{18}}$.
\end{solution}
\part $\displaystyle \left\{\frac56-\left[\frac34-\left(\frac12-\frac23\right)\right]\right\}:\left(-\frac16\right)+\frac25$.
\begin{solution}
Dall'interno: $-\frac16$, poi $\frac{11}{12}$, poi $-\frac1{12}$. Si ottiene $\frac12+\frac25=\boxed{\frac9{10}}$.
\end{solution}
\end{parts}
\end{esercizio}

\begin{esercizio}
\obiettivo{gestire potenze negative e frazioni complesse}
\question Calcola mostrando i passaggi.
\begin{parts}
\part $\displaystyle \left[\left(-\frac23\right)^{-2}\cdot\left(\frac23\right)^3\right]:\left(-\frac49\right)+\left(\frac12\right)^{-1}$.
\begin{solution}
$\frac94\cdot\frac8{27}=\frac23$; quindi $\frac23:(-\frac49)+2=-\frac32+2=\boxed{\frac12}$.
\end{solution}
\part $\displaystyle \frac{1-\frac34}{\frac56-\frac12}+\left(-\frac12\right)^{-2}-\left(\frac32\right)^0$.
\begin{solution}
$\frac{1/4}{1/3}+4-1=\frac34+3=\boxed{\frac{15}4}$.
\end{solution}
\end{parts}
\end{esercizio}

\newpage\sezione{Dal linguaggio comune all'espressione}

\begin{esercizio}
\obiettivo{tradurre istruzioni verbali}
\question Scrivi un'unica espressione che traduca le istruzioni, poi calcolane il valore.
\begin{parts}
\part Sottrai al quadrato della somma di $-3$ e $5$ il prodotto tra $-2$ e il valore assoluto di $-7$.
\begin{solution}
$(-3+5)^2-(-2)\cdot|-7|=4+14=\boxed{18}$.
\end{solution}
\part Dividi la differenza tra i tre quarti e i cinque sesti per l'opposto di un terzo; aggiungi poi il quadrato di meno un mezzo.
\begin{solution}
$\left(\frac34-\frac56\right):\left(-\frac13\right)+\left(-\frac12\right)^2=\frac14+\frac14=\boxed{\frac12}$.
\end{solution}
\end{parts}
\end{esercizio}

\begin{esercizio}
\obiettivo{modellizzare una successione di operazioni}
\question Per ciascun problema scrivi un'unica espressione e calcolane il valore.
\begin{parts}
\part Giulia possiede $756$ figurine. Regala a un'amica un settimo del numero iniziale, poi ne dà altre $156$ al fratello. Infine regala un terzo delle figurine rimaste. Quante figurine conserva?
\begin{solution}
$\left[756\left(1-\frac17\right)-156\right]\left(1-\frac13\right)=(648-156)\frac23=\boxed{328}$. Il secondo regalo frazionario riguarda il numero rimasto, non quello iniziale.
\end{solution}
\part Un serbatoio contiene $180$ litri d'acqua. Si utilizza un quarto dell'acqua iniziale, poi si aggiungono $18$ litri. Infine si utilizzano i due noni dell'acqua presente in quel momento. Quanti litri rimangono?
\begin{solution}
$\left[180\left(1-\frac14\right)+18\right]\left(1-\frac29\right)=153\cdot\frac79=\boxed{119}\text{ litri}$.
\end{solution}
\end{parts}
\end{esercizio}

\begin{esercizio}
\obiettivo{distinguere oggetti e intervalli}
\question Scrivi un'espressione e rispondi con l'unità di misura.
\begin{parts}
\part In una gara di $110$ m sono disposti $10$ ostacoli a intervalli uguali. Il primo è a $13{,}72$ m dalla partenza; l'ultimo è a $14{,}02$ m dal traguardo. Qual è la distanza fra due ostacoli consecutivi?
\begin{solution}
Fra $10$ ostacoli ci sono $9$ intervalli. $(110-13{,}72-14{,}02):9=\boxed{9{,}14\text{ m}}$.
\end{solution}
\part Su un percorso di $100$ m ci sono $10$ ostacoli. Il primo si trova a $13$ m dalla partenza e la distanza tra ostacoli consecutivi è $8{,}5$ m. Quanto misura il tratto dall'ultimo ostacolo al traguardo?
\begin{solution}
$100-13-9\cdot8{,}5=\boxed{10{,}5\text{ m}}$.
\end{solution}
\end{parts}
\end{esercizio}

\newpage\sezione{M.C.D., m.c.m. e coincidenze}

\begin{esercizio}
\obiettivo{calcolare M.C.D. e m.c.m.}
\question Scomponi in fattori primi e calcola M.C.D. e m.c.m. di ciascun gruppo.
\begin{parts}
\part $84,\ 126,\ 210$.
\begin{solution}
$84=2^2\cdot3\cdot7$, $126=2\cdot3^2\cdot7$, $210=2\cdot3\cdot5\cdot7$. M.C.D. $=\boxed{42}$; m.c.m. $=\boxed{1260}$.
\end{solution}
\part $72,\ 108,\ 150$.
\begin{solution}
$72=2^3\cdot3^2$, $108=2^2\cdot3^3$, $150=2\cdot3\cdot5^2$. M.C.D. $=\boxed6$; m.c.m. $=\boxed{5400}$. Nel M.C.D. si prendono i fattori comuni con esponente minimo, nel m.c.m. tutti i fattori con esponente massimo.
\end{solution}
\end{parts}
\end{esercizio}

\begin{esercizio}
\obiettivo{scegliere tra M.C.D. e m.c.m.}
\question Risolvi e spiega perché usi M.C.D. oppure m.c.m.
\begin{parts}
\part Si vogliono distribuire $84$ barrette e $126$ succhi nel maggior numero possibile di sacchetti tutti uguali, senza avanzi. Quanti sacchetti si preparano e che cosa contiene ciascuno?
\begin{solution}
Il numero dei sacchetti deve dividere entrambe le quantità e deve essere massimo: $\operatorname{MCD}(84,126)=42$. Ogni sacchetto contiene $2$ barrette e $3$ succhi.
\end{solution}
\part Due autobus partono insieme dal capolinea alle 8:10. Uno riparte ogni $18$ minuti e l'altro ogni $24$ minuti. A che ora ripartono insieme per la prima volta dopo le 8:10?
\begin{solution}
Il tempo cercato è il minimo multiplo positivo comune: $\operatorname{mcm}(18,24)=72$ minuti. Ripartono insieme alle \textbf{9:22}.
\end{solution}
\end{parts}
\end{esercizio}

\begin{esercizio}
\obiettivo{trovare coincidenze con tempi non interi}
\question Due segnali luminosi si accendono insieme all'istante iniziale. Il primo si riaccende ogni $5{,}5$ s, il secondo ogni $12{,}1$ s. Dopo quanti secondi si riaccenderanno insieme per la prima volta? Quanti intervalli completi saranno trascorsi per ciascun segnale? Mostra come riconduci il problema a numeri interi.
\begin{solution}
Misurando il tempo in decimi di secondo i periodi diventano $55$ e $121$. Poiché $55=5\cdot11$ e $121=11^2$, il m.c.m. è $605$ decimi, cioè $\boxed{60{,}5\text{ s}}$. Sono trascorsi $605:55=11$ intervalli del primo e $605:121=5$ del secondo. Il m.c.m. assicura che questa sia la prima coincidenza positiva.
\end{solution}
\end{esercizio}

\begin{esercizio}
\obiettivo{ricavare periodi razionali e confrontare andature}
\question Alice e Bob corrono sulla stessa pista, ad andatura costante. Alice compie $12$ giri in $8$ minuti, Bob $10$ giri in $6$ minuti. Partono insieme dalla linea di partenza.
\begin{parts}
\part Quanto dura un giro per ciascuno, in minuti e in secondi?
\part Dopo quanto tempo si troveranno entrambi per la prima volta di nuovo sulla linea di partenza? Quanti giri avrà completato ciascuno?
\part Chi è più veloce? Di quale percentuale la sua velocità supera quella dell'altro?
\end{parts}
\begin{solution}
Alice: $\frac8{12}=\frac23$ min $=40$ s per giro. Bob: $\frac6{10}=\frac35$ min $=36$ s per giro. $\operatorname{mcm}(40,36)=360$ s $=\boxed{6\text{ min}}$: Alice ha completato $9$ giri e Bob $10$. Bob è più veloce: nello stesso tempo percorre $\frac{10}{9}$ della distanza di Alice. L'incremento rispetto alla velocità di Alice è $(\frac{10}{9}-1)\cdot100\%=\boxed{\frac{100}{9}\%=11{,}\overline1\%}$. Non si cercano gli eventuali incontri in altri punti della pista.
\end{solution}
\end{esercizio}

\newpage\sezione{Rapporti, proporzioni e proporzionalità}

\begin{esercizio}
\obiettivo{risolvere e giustificare proporzioni}
\question Trova il termine incognito e verifica il risultato.
\begin{parts}
\part $\displaystyle \frac34:\frac56=x:\frac{10}{9}$.
\begin{solution}
$x=(\frac34\cdot\frac{10}{9}):\frac56=\boxed1$. I due rapporti valgono $\frac9{10}$.
\end{solution}
\part $\displaystyle 2{,}4:x=0{,}6:1{,}5$. Spiega inoltre perché da $\frac ab=\frac cd$, con $b,d\ne0$, segue $ad=bc$, usando la proprietà invariantiva e un denominatore comune.
\begin{solution}
$x=(2{,}4\cdot1{,}5):0{,}6=\boxed6$. I rapporti valgono $0{,}4$. Per la proprietà invariantiva $\frac ab=\frac{ad}{bd}$ e $\frac cd=\frac{bc}{bd}$; frazioni uguali con lo stesso denominatore non nullo hanno numeratori uguali, quindi $ad=bc$.
\end{solution}
\end{parts}
\end{esercizio}

\begin{esercizio}
\obiettivo{riconoscere proporzionalità diretta e inversa}
\question Risolvi e indica il tipo di proporzionalità e la quantità che resta costante.
\begin{parts}
\part Un tessuto costa $18$ euro per $2{,}4$ m. Il prezzo al metro è costante. Quanto costano $3{,}6$ m? Quanti metri si acquistano con $30$ euro?
\begin{solution}
Proporzionalità diretta; rapporto prezzo/lunghezza $=18:2{,}4=7{,}5$ euro/m. $3{,}6$ m costano $\boxed{27}$ euro; con $30$ euro si acquistano $\boxed4$ m.
\end{solution}
\part Sei rubinetti identici, aperti insieme, riempiono una vasca in $15$ ore. Quanto impiegano $10$ rubinetti identici? Quanti rubinetti occorrono per riempirla in $5$ ore? Supponi costanti la portata di ogni rubinetto e il volume della vasca.
\begin{solution}
Proporzionalità inversa; prodotto numero di rubinetti per tempo $=6\cdot15=90$. Con $10$ rubinetti servono $\boxed9$ ore; in $5$ ore occorrono $\boxed{18}$ rubinetti.
\end{solution}
\end{parts}
\end{esercizio}

\begin{esercizio}
\obiettivo{riconoscere modelli in tabelle}
\question Completa le tabelle e specifica se la proporzionalità è diretta o inversa. Scrivi la relazione fra $x$ e $y$.
\begin{parts}
\part $\begin{array}{c|rrrr}x&2&5&8&12\\\hline y&3&\square&12&\square\end{array}$
\begin{solution}
Diretta: $y/x=\frac32$, dunque $y=\frac32x$. Mancano $\boxed{7{,}5}$ e $\boxed{18}$.
\end{solution}
\part $\begin{array}{c|rrrr}x&2&3&6&12\\\hline y&18&\square&6&\square\end{array}$
\begin{solution}
Inversa: $xy=36$, dunque $y=36/x$. Mancano $\boxed{12}$ e $\boxed3$.
\end{solution}
\end{parts}
\end{esercizio}

\begin{esercizio}
\obiettivo{usare scale e rapporti}
\question Risolvi mostrando l'impostazione.
\begin{parts}
\part Su una carta in scala $1:25\,000$ un sentiero misura $7{,}2$ cm. Quanto è lungo nella realtà, in km? Quanto misura sulla carta un tratto reale di $750$ m?
\begin{solution}
$7{,}2\cdot25\,000=180\,000$ cm $=\boxed{1{,}8\text{ km}}$. $750$ m $=75\,000$ cm: sulla carta $75\,000:25\,000=\boxed{3\text{ cm}}$.
\end{solution}
\part Per una bevanda si mescolano sciroppo e acqua nel rapporto $2:7$. Si vogliono preparare $1{,}8$ litri di bevanda, assumendo che i volumi si sommino. Quanto sciroppo e quanta acqua occorrono?
\begin{solution}
Le parti sono $2+7=9$; ciascuna vale $1{,}8:9=0{,}2$ litri. Sciroppo $\boxed{0{,}4\text{ L}}$, acqua $\boxed{1{,}4\text{ L}}$.
\end{solution}
\end{parts}
\end{esercizio}

\newpage\sezione{Percentuali e problemi a più passaggi}

\begin{esercizio}
\obiettivo{calcolare parte, percentuale e totale}
\question Risolvi mostrando i calcoli.
\begin{parts}
\part In una scuola con $240$ studenti il $35\%$ partecipa a un torneo. Quanti sono i partecipanti? Fra questi, $21$ disputano la finale: quale percentuale dei partecipanti rappresentano?
\begin{solution}
$240\cdot\frac{35}{100}=\boxed{84}$ partecipanti; $\frac{21}{84}\cdot100\%=\boxed{25\%}$.
\end{solution}
\part Il $30\%$ degli iscritti a un corso corrisponde a $18$ persone. Quanti sono gli iscritti? Quale percentuale del totale rappresentano $27$ persone?
\begin{solution}
Totale $18:\frac{30}{100}=\boxed{60}$. Le $27$ persone sono $\frac{27}{60}\cdot100\%=\boxed{45\%}$.
\end{solution}
\end{parts}
\end{esercizio}

\begin{esercizio}
\obiettivo{gestire variazioni percentuali successive}
\question Scrivi un'espressione e rispondi.
\begin{parts}
\part Una felpa costa $80$ euro. Il prezzo aumenta del $20\%$ e successivamente diminuisce del $20\%$ rispetto al prezzo aumentato. Qual è il prezzo finale? Qual è la variazione percentuale complessiva rispetto agli $80$ euro iniziali?
\begin{solution}
$80(1+\frac{20}{100})(1-\frac{20}{100})=80\cdot1{,}2\cdot0{,}8=\boxed{76{,}80\text{ euro}}$. La diminuzione è $3{,}20$ euro, cioè $\boxed{4\%}$ del prezzo iniziale. Le due percentuali hanno basi diverse.
\end{solution}
\part Dopo uno sconto del $25\%$ un articolo costa $63$ euro. Qual era il prezzo iniziale? Di quale percentuale dovrebbe aumentare il prezzo scontato per tornare a quello iniziale?
\begin{solution}
$63:(1-\frac{25}{100})=\boxed{84\text{ euro}}$. Per risalire da $63$ a $84$, l'aumento è $\frac{21}{63}\cdot100\%=\boxed{\frac{100}{3}\%=33{,}\overline3\%}$.
\end{solution}
\end{parts}
\end{esercizio}

\begin{esercizio}
\obiettivo{calcolare percentuali riferite a gruppi diversi}
\question In un istituto il $40\%$ degli studenti è costituito da ragazze e il restante $60\%$ da ragazzi. Vuole proseguire gli studi il $35\%$ delle ragazze e il $30\%$ dei ragazzi.
\begin{parts}
\part Quale percentuale di tutti gli studenti vuole proseguire gli studi?
\part Le ragazze che vogliono proseguire sono $28$. Quanti sono in tutto gli studenti? Quanti ragazzi vogliono proseguire?
\end{parts}
\begin{solution}
Su $100$ studenti, $40\cdot0{,}35=14$ ragazze e $60\cdot0{,}30=18$ ragazzi vogliono proseguire: totale $\boxed{32\%}$. Le ragazze totali sono $28:0{,}35=80$, dunque gli studenti sono $80:0{,}40=\boxed{200}$. I ragazzi sono $120$, dei quali $\boxed{36}$ vogliono proseguire. Non si fa la media semplice fra $35\%$ e $30\%$.
\end{solution}
\end{esercizio}

\begin{esercizio}
\obiettivo{individuare la base di una percentuale}
\question In un gruppo sportivo di $80$ iscritti, il $60\%$ pratica atletica. Fra chi pratica atletica, il $25\%$ partecipa a una gara. Quanti iscritti partecipano alla gara? Quale percentuale degli $80$ iscritti rappresentano? Uno studente propone $60\%+25\%=85\%$: spiega l'errore.
\begin{solution}
Gli atleti sono $80\cdot0{,}60=48$; i partecipanti sono $48\cdot0{,}25=\boxed{12}$. Rispetto al gruppo iniziale sono $\frac{12}{80}\cdot100\%=\boxed{15\%}$. Il $25\%$ si riferisce solo al sottogruppo del $60\%$: le percentuali si moltiplicano, $0{,}60\cdot0{,}25=0{,}15$, e non si sommano.
\end{solution}
\end{esercizio}
\end{esercizi}
\finefila
\end{document}









